\begin{explanationbox}
\textbf{\textcolor{CExplanation}{先想一个生活画面}}

把开口向上的抛物线想成一个\textbf{U 形滑梯}，最低处是谷底。现在只允许你在滑梯上的\textbf{一段}活动（这段就是区间）：谷底在活动范围内，最低点自然是谷底；谷底在外面，你根本走不到，只能走到\textbf{靠它最近的一端}。找最高点则要去\textbf{离谷底最远的一端}。如果倒过来就是开口向下。

\medskip
\textbf{\textcolor{CExplanation}{用具体数字看懂“区间动”}}

固定 $f(x)=(x-2)^2$，轴一直是 $x=2$，让长度为 $2$ 的区间 $[t,t+2]$ 从左往右移动：
\begin{center}
\begin{tikzpicture}[x=0.78cm,y=0.26cm,>=stealth]
  \draw[->,gray] (-1.5,0)--(6.3,0) node[right] {$x$};
  \draw[gray!60,dashed] (2,-1.35)--(2,17) node[above] {$x=2$};
  \draw[thick,teal!70!black,domain=-1.3:5.8,samples=85,smooth] plot (\x,{(\x-2)^2});
  \draw[line width=2.5pt,orange!90!black] (-1,-1.15)--(1,-1.15);
  \draw[line width=2.5pt,blue!70!black] (1,-2.1)--(3,-2.1);
  \draw[line width=2.5pt,red!75!black] (3,-3.05)--(5,-3.05);
  \foreach \x in {-1,1,2,3,5} {\fill[gray!70] (\x,0) circle (1.3pt);}
  \node[left,orange!85!black] at (-1,-1.15) {左边};
  \node[right,blue!70!black] at (3,-2.1) {跨轴};
  \node[right,red!75!black] at (5,-3.05) {右边};
  \node[below] at (2,-0.05) {$2$};
\end{tikzpicture}

\end{center}
\begin{itemize}[leftmargin=*,itemsep=0.35em]
  \item 区间 $[-1,1]$：不包含 $2$，离轴近的是 $1$，最小值 $f(1)=1$；远的是 $-1$，最大值 $f(-1)=9$。
  \item 区间 $[1,3]$：包含 $2$，最小值 $f(2)=0$；两端到轴距离同为 $1$，\textbf{最大值两端都能取得}，等于 $1$。
  \item 区间 $[3,5]$：不包含 $2$，最小值 $f(3)=1$，最大值 $f(5)=9$。
\end{itemize}

\textbf{\textcolor{CExplanation}{为什么比较“横坐标距离”就够？}}

先配方：$f(x)=a(x-h)^2+k$。纵坐标只由\textbf{$x$ 到 $h$ 的距离的平方}决定。比如 $f(x)=(x-2)^2$，$x=0$ 与 $x=4$ 到轴的距离都为 $2$，所以函数值都为 $4$。开口向上时，离轴越远越高；开口向下时，离轴越远越低。

\medskip
\textbf{\textcolor{CExplanation}{眼前一亮：为什么“中点”会让端点换班？}}

只比较左右端点谁更远：\textbf{若区间中点在轴左侧，左端点更远；若中点在轴右侧，右端点更远；中点恰在轴上时两端一样高。}
\[
\boxed{\text{换班位置：}\ \frac{m+n}{2}=h}
\]
\begin{center}
\begin{tikzpicture}[x=1.0cm,y=0.62cm,>=stealth]
 \draw[->,gray] (-2.8,0)--(3.0,0) node[right] {$x$};
 \draw[dashed,gray!65] (0,-0.4)--(0,2.6);
 \draw[teal!70!black,very thick,domain=-2.55:2.55,samples=70] plot (\x,{0.32*\x*\x});
 \draw[line width=2.5pt,orange!85!black] (-1.8,-0.28)--(1.8,-0.28);
 \fill[red!70!black] (-1.8,1.0368) circle (2.5pt) (1.8,1.0368) circle (2.5pt);
 \node[below] at (-1.8,-0.35) {$m$};
 \node[below] at (0,-0.35) {$h=0$};
 \node[below] at (1.8,-0.35) {$n$};
 \node[above] at (0,2.55) {左右等远 $\Rightarrow$ 两端等高};
\end{tikzpicture}

\end{center}

\textbf{\textcolor{CExplanation}{最快做题顺序}}

\textbf{第一眼}看开口和固定的轴；\textbf{第二眼}看区间能否碰到轴；\textbf{第三眼}比较端点到轴的距离。遇到参数，重点解\textbf{端点碰轴}和\textbf{区间中点碰轴}这两类等式。\textit{本结论是“轴固定、区间动”；F031 处理的是“区间固定、轴动”，不要混淆。}
\end{explanationbox}
